% Apéndice: composición afín práctica: rotación y escala
\section{Composición afín práctica: rotación y escala}\label{cap:16}\label{ap:composicion-afin}

\subsection{Rotación alrededor de un centro}

Sea \(c=(c_x,c_y)^T\) y sea \(R\) una matriz de rotación. Trasladar \(c\) al origen, rotar
y trasladar de regreso produce
\[
F_1(x)=c+R(x-c)=Rx+(I-R)c.
\]
En coordenadas homogéneas:
\[
\boxed{
H_1=
\left(
\begin{array}{c|c}
R&(I-R)c\\ \hline
0\quad0&1
\end{array}
\right)
}.
\]
La escritura por bloques significa que \(R\) ocupa las primeras dos filas y columnas, mientras
que el vector \((I-R)c\) ocupa la última columna.

Para \(\theta=45^\circ\):
\[
R=
\begin{pmatrix}
r&-r\\
r&r
\end{pmatrix},
\qquad
r=\frac{\sqrt2}{2}.
\]

\begin{atencion}[Error frecuente con \(45^\circ\)]
Es común escribir en algún pasaje \(\sqrt3/2\). Para \(45^\circ\), el valor correcto es
\[
\sin45^\circ=\cos45^\circ=\frac{\sqrt2}{2}.
\]
\(\sqrt3/2\) corresponde al coseno de \(30^\circ\) o al seno de \(60^\circ\).
\end{atencion}

\subsection{Escala uniforme alrededor de otro centro}

Sea \(d=(d_x,d_y)^T\) y sea \(k\) el factor de escala:
\[
F_2(y)=d+k(y-d)=ky+(1-k)d.
\]
Su matriz homogénea es
\[
\boxed{
H_2=
\begin{pmatrix}
k&0&(1-k)d_x\\
0&k&(1-k)d_y\\
0&0&1
\end{pmatrix}
}.
\]
Con \(k=0.5\), cada distancia al centro \(d\) queda reducida a la mitad.

\subsection{Composición y orden}

La rotación se aplica primero y la escala después:
\[
F=F_2\circ F_1.
\]
Por eso, con vectores columna,
\[
\boxed{H=H_2H_1}.
\]
La matriz situada a la derecha actúa primero. Al desarrollar por bloques:
\[
F(x)=kRx+\left[k(I-R)c+(1-k)d\right].
\]
Así,
\[
A=kR,
\qquad
t=k(I-R)c+(1-k)d.
\]

\subsection{Ejemplo numérico completo}

Se eligen
\[
c=(3,1),\qquad d=(-1,2),\qquad
k=0.5,\qquad \theta=45^\circ.
\]
La rotación alrededor de \(c\) queda
\[
H_1=
\begin{pmatrix}
0.70710678&-0.70710678&1.58578644\\
0.70710678& 0.70710678&-1.82842712\\
0&0&1
\end{pmatrix}.
\]
La escala respecto de \(d\) es
\[
H_2=
\begin{pmatrix}
0.5&0&-0.5\\
0&0.5&1\\
0&0&1
\end{pmatrix}.
\]
El producto es
\[
\boxed{
H=H_2H_1=
\begin{pmatrix}
0.35355339&-0.35355339&0.29289322\\
0.35355339& 0.35355339&0.08578644\\
0&0&1
\end{pmatrix}
}.
\]
Por tanto:
\[
A=
\begin{pmatrix}
0.35355339&-0.35355339\\
0.35355339& 0.35355339
\end{pmatrix},
\qquad
t=
\begin{pmatrix}
0.29289322\\
0.08578644
\end{pmatrix}.
\]

\subsection{Punto fijo}

Un punto \(p\) es fijo si
\[
Ap+t=p.
\]
Por lo tanto:
\[
(I-A)p=t,
\qquad
\boxed{p=(I-A)^{-1}t}.
\]
En este ejemplo:
\[
\boxed{
p_\ast=
\begin{pmatrix}
1-\frac{\sqrt2}{2}\\[1mm]
1-\frac{\sqrt2}{2}
\end{pmatrix}
\approx
\begin{pmatrix}
0.29289322\\
0.29289322
\end{pmatrix}
}.
\]
La verificación directa es
\[
Hp_\ast^{\,h}=p_\ast^{\,h},
\qquad
p_\ast^{\,h}=
\begin{pmatrix}p_{\ast x}\\p_{\ast y}\\1\end{pmatrix}.
\]

\subsection{Autovalores y coordenadas homogéneas}

Como \(H\) es triangular por bloques, sus autovalores son los dos autovalores de \(A=0.5R\)
y el autovalor homogéneo \(1\):
\[
\lambda_{1,2}=0.5e^{\pm i\pi/4},
\qquad
\lambda_3=1.
\]
Los dos autovalores complejos describen conjuntamente la contracción y el giro de la parte
lineal. El autovector real asociado a \(\lambda_3=1\) puede escribirse
\[
v_3=
\begin{pmatrix}p_{\ast x}\\p_{\ast y}\\1\end{pmatrix},
\]
y representa el punto fijo.

\begin{atencion}[Cómo interpretar el resultado de una calculadora online]
Los autovectores se devuelven con una normalización arbitraria. Si la calculadora entrega
\[
v=(u,v,w)^T,\qquad w\neq0,
\]
el punto cartesiano correspondiente es \((u/w,v/w)\). Si \(w=0\), el vector representa una
dirección o punto ideal, no un punto fijo finito. Una matriz homogénea afín siempre posee
\(1\) como autovalor, pero solo hay punto fijo finito cuando existe un autovector asociado
con \(w\neq0\).
\end{atencion}

\subsection{Implementación en Excel}

\begin{enumerate}
  \item Ingresar \(H_1\) en un bloque \(3\times3\), por ejemplo
  \texttt{B2:D4}.
  \item Ingresar \(H_2\) en \texttt{F2:H4}.
  \item Seleccionar un bloque \(3\times3\) para el resultado e introducir
  \[
  \texttt{=MMULT(F2:H4,B2:D4)}.
  \]
  En versiones antiguas de Excel se confirma con
  \texttt{Ctrl+Shift+Enter}; en versiones con matrices dinámicas basta \texttt{Enter}.
  \item Para transformar un punto homogéneo almacenado como columna
  \((x,y,1)^T\), usar
  \[
  \texttt{=MMULT(matriz\_compuesta,rango\_del\_punto)}.
  \]
  \item Para obtener el punto fijo sin recurrir a una web, construir \(I-A\) y calcular
  \[
  \texttt{=MMULT(MINVERSE(rango\_I\_menos\_A),rango\_t)}.
  \]
  El nombre de \texttt{MINVERSE} puede variar según el idioma de Excel.
\end{enumerate}

\begin{idea}[Comprobación mínima]
Transformar todos los vértices de la figura, verificar uno manualmente y comprobar que
\[
Hp_\ast^{\,h}=p_\ast^{\,h}.
\]
En el programa gráfico, los comandos relevantes son \texttt{ROTATE} y
\texttt{SCALE}. El estilo o tamaño del punto puede aumentarse para hacer visibles centros y
vértices.
\end{idea}

